\section{Introduction}
\label{sec:introduction}

Crews and rovers rehearse on Earth before they work on the Moon or Mars, and the value of a
rehearsal depends on how well the Earth site stands in for the target. Connecting Ridge, one of
the candidate Artemis III landing regions that NASA announced near the lunar south pole \citep{nasa2022artemis}, has elevation models at 5 m per pixel built from laser altimetry alone \citep{barker2021}. The
practical question is which Earth surfaces are worth visiting to learn about such a site, and on
what evidence. Quantitative studies show how hard the comparison is: even three lava facies at
one Icelandic flow field overlap by 88\% to 89\% in their radar polarization ratio
\citep{tolometti2022holuhraun}.

Roughness statistics are the natural quantitative starting point. Roughness has been reported at stated baselines from terrain data on Earth, the Moon and Mars \citep{shepard2001,rosenburg2011,kreslavsky2000}, and the structure function of elevation and its Hurst exponent can be computed with shared definitions on all three bodies, while each product keeps its own instrument, processing and support. They are second order statistics, however. They summarise how height
differences grow with lag and do not determine whether a surface is cratered, ridged or covered in flows, so two windows can be close in every statistic and still differ in the morphology that matters for a traverse. Learned representations can carry morphology. Recent lunar foundation models are trained on lunar data \citep{gironamata2026,fraccaro2026}, and a learned terrain analog retrieval works between two terrestrial regions \citep{liu2026}; a learned similarity across bodies is only useful if it can be checked against physics.

We take up both routes for one lunar target, under one set of declared rules: every distance names
its observable and its supported lags, missing evidence is never read as zero, and every number in
this paper is generated from a recorded artifact. The fingerprint is computed by one Python
implementation offline and by a TypeScript port that runs in a browser worker, and the two are held
to stated numerical tolerances (Section~\ref{sec:benchmarks}). The contributions are as follows.
First, a deterministic fingerprint of relief and texture compares the Connecting Ridge target
window with \ComparedCount{} eligible Earth windows and returns a shortlist of
\ShortlistCount{} candidate window, with the closest relief distance \BestReliefDistance{} against a median of \ReliefDistanceMedian{} (Section~\ref{sec:results-retrieval}); the same comparison for the Jezero and Arcadia windows returns shortlists of \JezeroShortlistSize{} and \ArcadiaShortlistSize{} windows whose distances lie below the median internal distance of each Martian site (Section~\ref{sec:results-mars}). Second, a radar availability audit decodes the
Mini-RF mosaic over the site and a UAVSAR scene over Holuhraun, and records why a radar comparison
could not be qualified within the study (Section~\ref{sec:results-radar}). Third, a terrain encoder for Moon, Mars and Earth elevation with a physics head and a gradient reversal body head is held to acceptance criteria fixed before training; trained at a reduced scope, its probe results are reported against those criteria with an untrained copy of the network as a baseline, and the morphology probe and the image baselines are not available (Section~\ref{sec:results-ml}). Fourth, a benchmark section reports a validation sweep over the adversarial weight and training length, training throughput on a CPU server and a laptop accelerator, and the numerical parity of the browser implementations (Section~\ref{sec:benchmarks}).
